\chapter{Conclusion}
\label{chap:conclusion}

This thesis addressed the problem of synthesizing plausible three-dimensional human--scene interactions without task-specific training. Given a natural-language instruction and a single posed view of a reconstructed indoor scene, the goal was to place a single SMPL-X body into the scene so that it performs the described interaction, contacting the intended surfaces while remaining physically plausible. The difficulty is that the interaction knowledge needed to solve this lives in the space of images and language, where large pretrained models are strong, while the answer must be expressed in the metric three-dimensional space of the scene, where those models do not operate. Bridging that gap without training on paired human--scene data was the central challenge of this work.

The approach rests on a single idea: rather than using foundation models to produce the interacting body directly, the method uses them to \emph{localize} where the interaction touches the scene, and then grounds those localized contacts on the reconstructed geometry through the calibrated camera. A human performing the interaction is generated into the posed view, the body-part-specific contact regions are marked on that view and projected to the mesh as explicit optimization targets, and a scene-aware optimization refines a body to meet those targets while avoiding penetration. Because the contacts are localized in the same view whose geometry is known, image-space reasoning is converted into metric three-dimensional targets exactly, so that generation supplies the interaction knowledge while camera geometry and optimization turn it into a physically correct body.

Making these localized contacts reliable is the responsibility of the method's central contribution, an agentic contact-estimation procedure that does not trust a single generated contact annotation but instead verifies and corrects it in a closed loop, so that an unreliable single generation is turned into a dependable set of precise, per-body-part contact regions. This procedure forms the core of a method comprising five modules. Within it, a Scene Interaction Graph provides only a coarse, symbolic account of which body parts contact the scene, deliberately kept as lightweight scaffolding, while the precise surface regions are recovered visually and verified downstream. The remainder of this chapter states how the results answer the research questions of \cref{sec:intro-rq}, summarizes the contributions, and discusses the limitations and the directions they open.

\section{Answers to the Research Questions}
\label{sec:conclusion-rq}

The evaluation of \cref{chap:experiments} was structured around the three sub-questions posed in \cref{sec:intro-rq}. Each can now be answered from the evidence gathered there, without reproducing the numbers.

\paragraph{Localization.} The first sub-question asked whether an image generation model can localize fine-grained, body-part-specific contact regions on a scene surface reliably enough that they can be projected into 3D and used as optimization targets. The evidence indicates that it can. The qualitative study of the contact-estimation stage showed the localized regions placed on the correct object surfaces, and the low mean contact distance attained by the full method confirmed that the resulting 3D targets are accurate enough to draw the body into contact with the intended regions. Localization is therefore feasible in the sense the question requires. The qualification is that it depends on the contacted region being visually identifiable: on objects with many near-identical parts, such as the rungs of a ladder, the correct region can be under-determined by the image, which is the source of one of the observed failure cases.

\paragraph{Verification.} The second sub-question asked whether a closed-loop, agentic verify-and-correct procedure produces more reliable contact regions than a single generation. The ablation was designed to isolate exactly this comparison, holding the overlay model, the optimization, and the initialization fixed and varying only whether the contacts are verified and corrected before use. Adding the agentic loop reduced the mean contact distance substantially while leaving penetration essentially unchanged, so the improvement fell specifically on the accuracy of the scene-side contact rather than on any other property of the result. This is the central claim of the thesis, and the evidence supports it: verifying and correcting the contacts, rather than accepting a single generation, is what makes the localized contact dependable, and it is the agentic loop that turns an unreliable single generation into a usable optimization target.

\paragraph{Grounding.} The third sub-question asked whether grounding the interaction on the reconstructed scene geometry, through these localized contacts, yields bodies that are more physically plausible than those obtained by fitting to image evidence alone. The comparison with the GVHMR initialization answers this directly. The initialization is a body recovered from the human frame without knowledge of the scene, and it attained both the worst contact distance and the highest penetration of any setting; grounding that body on the scene through the localized contacts reduced both. The comparison with the PhySIC baseline pointed the same way, with the explicitly grounded contacts yielding body parts placed more precisely on the intended surfaces than a method that recovers contact implicitly. Grounding the interaction on the reconstructed geometry therefore delivers the physical plausibility that image evidence alone does not, which was the premise motivating the whole approach.

Taken together, the three answers trace the logic of the method from end to end: contact can be localized in image space, a verification loop makes that localization reliable, and grounding the verified contacts on the reconstructed scene produces a body that is physically plausible in the scene. This is a positive answer to the central research question of \cref{sec:intro-rq}: a zero-shot method built only from pretrained foundation models and explicit camera geometry can recover precise, verified, per-body-part contact regions on a reconstructed scene and use them to synthesize a plausible static human--scene interaction in the scene's coordinate frame.

\section{Contributions}
\label{sec:conclusion-contributions}

The contributions of this thesis follow from the answers above and can be stated concisely.

\begin{enumerate}
  \item \textbf{Agentic contact estimation.} The central contribution is a procedure that localizes precise, per-body-part contact regions on a scene surface and makes them reliable through a closed verify-and-correct loop. An image generation model proposes body-part-specific contact regions, which are composited onto the untouched posed view to keep them in exact correspondence with the calibrated camera, and a vision-language evaluator inspects each result against the intended interaction and issues corrections until the contacts are accepted. Because the regions are localized visually rather than inferred from the identity of a named element, and because they are verified before use, the procedure resolves contact on small and large functional surfaces alike and turns an unreliable single generation into a dependable optimization target. The ablation showed this loop to be the component responsible for the accuracy of the final scene-side contact.

  \item \textbf{A zero-shot method for scene-grounded interaction synthesis.} The thesis composes pretrained foundation models with explicit camera geometry and a scene-aware optimization into a unified method that synthesizes a static human--scene interaction from a language instruction, a single posed view, and a reconstructed mesh, without any interaction-specific training. The method lifts the visually localized contacts to the mesh through the calibrated camera and refines a body to meet them while avoiding penetration and staying close to a monocular initialization, converting image-space interaction knowledge into a body that is physically correct in the scene coordinate frame.

  \item \textbf{An evaluation methodology for scene-side contact.} Assessing whether a body contacts the correct scene regions is not straightforward in this setting, since no contact ground truth exists for these interactions and standard geometric metrics alone do not capture whether the depicted action is correct. The thesis addresses this with two evaluation components. The first is a manual scene-side contact annotation that mirrors the method's own contact representation, painting the region each body part should contact directly onto the posed view, which provides a per-body-part reference against which the method's localized contacts can be measured on the mesh. The second is a vision-language rubric evaluation that scores each interaction along target-object, action, contact, and physical-plausibility criteria with explicit caps for characteristic failures, serving as a scalable proxy for perceptual judgment alongside the geometric measures. Together these make the scene-side quality of a synthesized interaction measurable, and both are applied identically to the proposed method, its ablations, and the external baseline.
\end{enumerate}

\section{Limitations and Future Work}
\label{sec:conclusion-future}

The method has clear limitations, and the analysis of \cref{sec:exp-failure} identified two in particular: contact can be localized on the wrong part of an object when the object has many near-identical parts, such as the rungs of a ladder, and the result can be incorrect when the monocular initialization returns a pose far from the interaction, as in the case of a person lying on a sofa, from which the pose prior prevents the optimization from recovering. Beyond these, the method depends on external foundation models and was compared against a single external baseline. These observations motivate several directions for future work.

\begin{enumerate}
  \item \textbf{Open-source model ablation.} The method is model-agnostic by design: its language, image-generation, and vision-language stages are swappable components accessed through their interfaces, and the current system fills them with frontier, largely closed models. A natural next step is to replace these components with open-source counterparts, for instance an open inpainting and image-generation model in place of the proprietary generators and an open vision-language model in place of the proprietary evaluator and verifier, and to measure the resulting change in performance. Because the components are interchangeable, this is a controlled ablation rather than a redesign, and it would quantify how much each frontier model contributes to the final result and whether a fully open, self-hostable implementation is viable.

  \item \textbf{User study.} The vision-language rubric evaluation is used as a scalable proxy for human perceptual judgment, and its limitations were stated plainly: a vision-language judge may miss subtle penetration, be influenced by rendering, or apply the rubric inconsistently. A controlled user study, in which people rate the synthesized interactions along the same criteria, would establish how well the automated rubric tracks human judgment and would provide a perceptual reference against which the geometric metrics and the proxy can be calibrated.

  \item \textbf{Repeated sampling and variance.} Each interaction was synthesized once, and the results are reported as means over the evaluation set. Several stages of the method are stochastic, however, including the frame generation, the candidate selection, and the contact overlays, so a single run per interaction leaves the run-to-run variability of the method unmeasured. Synthesizing each interaction several times would report the metrics with their variance, establishing how much of the observed improvement is stable rather than particular to one sample. It would also extend what the diversity measures can show. As reported here they are computed across distinct instructions and serve to confirm that the method does not collapse to a narrow set of poses, which is the use made of them in \cref{sec:exp-quantitative}. Sampling repeatedly from a single instruction would additionally reveal how varied the solutions are for one interaction, a property the present protocol cannot observe.
  
  \item \textbf{Comparison with additional baselines.} The empirical comparison was made against a single external baseline, PhySIC, chosen as the closest method that produces an SMPL-X body in metric contact with scene geometry from a single view. The zero-shot methods positioned against the proposed approach in \cref{chap:related-work}, GenZI and FunHSI, were discussed conceptually but not evaluated empirically, and trained human--scene interaction methods were excluded as out of the zero-shot setting. For FunHSI, no public implementation is available. For GenZI, the obstacle is its generator: the method's own stated limitation is the inpainting capability of latent diffusion models, and that limitation is what the present setting exercises, since Stable Diffusion~2 inpaints poorly on ScanNet++ views, which are reconstructions rather than photographs. Running it unmodified would therefore measure its generator's difficulty with these inputs rather than its grounding strategy, which is the property the positioning of \cref{chap:related-work} actually claims to improve upon. A meaningful comparison would first require porting its dynamic masking to a current image generator, so that the two methods are separated by how they resolve the scene side of the contact rather than by the age of the models they were built on. Doing so, together with a comparison against FunHSI once an implementation is released and against representative trained baselines, would situate the method more completely and test the conceptual advantages claimed in the positioning against measured results.

  \item \textbf{Dynamic and motion extension.} The method synthesizes a single static interaction state, a deliberate choice that isolates the contact problem from the additional difficulty of temporal coherence. A natural extension is to move from a static state to motion, generating the trajectory that leads into or out of the interaction while keeping the contact grounded on the scene throughout. The verified scene-side contacts recovered here could serve as anchors for such a motion, defining the moments of contact that a generated trajectory must pass through.
\end{enumerate}

These directions share a common thread: the core of the method, the recovery of precise, verified, per-body-part contact regions on a reconstructed scene, is independent of the particular models that implement it, of the static setting in which it was evaluated, and of the specific baseline it was compared against. Each direction extends the same idea, whether by changing the models that realize it, validating the way it is measured, broadening what it is compared against, or carrying its contacts forward in time.

\section*{Closing Remarks}
 
This thesis set out from a simple observation: reconstructed scenes are empty, and populating them convincingly depends less on rendering a plausible-looking person than on establishing where, precisely, the body meets the scene. The approach taken here treats that contact as the quantity to solve for, and uses foundation models not to draw the final interaction but to localize and verify the contact regions that make it physically correct. That contact can be recovered this way, from a single view and a language instruction, without any interaction-specific training, and grounded reliably enough on reconstructed geometry to place a plausible body in the scene, is the central finding of this work. The specific models it relies on will be superseded, but the idea that precise, verified scene-side contact is what turns image-space interaction knowledge into a physically grounded body is what this thesis contributes, and it is one that the directions above can carry considerably further.
 
